% {¿Por qué es importante el proyecto? ¿Por qué esas tecnologías?
%  Incluir referencias.}
% [Mínimo media cuartilla, máximo dos]
El rendimiento de un equipo profesional de baloncesto no depende únicamente de las habilidades individuales
de sus jugadoras, sino también de las decisiones estratégicas relacionadas con las alineaciones y los
emparejamientos en cada encuentro. La selección adecuada de combinaciones en cancha puede maximizar
ventajas tácticas y mejorar la eficiencia colectiva; sin embargo, estas decisiones suelen basarse principalmente
en la experiencia del cuerpo técnico y en estadísticas tradicionales, lo que puede limitar el aprovechamiento
integral de los datos disponibles \cite{berri2006wages}.

El presente proyecto propone el desarrollo de un modelo de optimización de alineaciones basado en el análisis
de emparejamientos y configuraciones de cuadro en la WNBA, usando técnicas de aprendizaje automático,
ciencia de datos y analítica deportiva avanzada. La finalidad es generar una herramienta de apoyo a la toma de
decisiones que permita fundamentar estratégicamente las rotaciones y combinaciones de jugadoras,
contribuyendo a que los resultados dependan en mayor medida del desempeño y potencial de las atletas.
En términos de originalidad, el proyecto integra análisis de emparejamientos individuales con técnicas de
optimización y modelado predictivo, superando enfoques meramente descriptivos. Los beneficiarios
potenciales incluyen cuerpos técnicos, analistas deportivos y organizaciones interesadas en mejorar sus
procesos de decisión con base en evidencia cuantitativa. En este aspecto servirá como base para otros proyectos
y la información que genere sustentará la base a nuevas herramientas.
La viabilidad del proyecto se sustenta en la disponibilidad de datos históricos, herramientas computacionales
accesibles y metodologías consolidadas en ciencia de datos, lo que garantiza su desarrollo dentro de los tiempos
y recursos establecidos.
Los modelos matemáticos clásicos se basan únicamente en supuestos predefinidos o métricas agregadas;
mientras que los métodos no supervisados permiten capturar relaciones multivariadas complejas y dependencias
implícitas entre jugadoras, lo que resulta especialmente relevante en un deporte altamente interactivo como el
baloncesto \cite{yamada2024offensivelineupanalysisbasketball}.
La complejidad del proyecto radica en la integración de diferentes niveles de análisis dentro de un mismo
modelo: por un lado, la caracterización avanzada de perfiles de jugadoras a partir de múltiples indicadores de
rendimiento; y por otro, la evaluación de configuraciones colectivas y emparejamientos entre equipos. Este proceso implica la combinación de técnicas de análisis exploratorio, modelado estadístico y optimización, con
el fin de identificar patrones de sinergia entre jugadoras y estimar su impacto potencial en distintas alineaciones.
De esta manera, el trabajo no solo requiere la implementación de algoritmos de aprendizaje automático, sino
también el diseño de un marco analítico que permita interpretar los resultados en términos tácticos y
estratégicos.
La propuesta representa una mejora respecto a enfoques existentes centrados únicamente en el análisis
descriptivo o en la predicción de resultados, ya que introduce un componente orientado a identificar
configuraciones óptimas de jugadoras en función del contexto competitivo y del rival. La literatura en analítica
deportiva ha destacado la importancia de modelar interacciones colectivas y estructuras de juego para
comprender el rendimiento real de los equipos, más allá de estadísticas individuales aisladas \cite{skinner2010price}. Bajo esta
perspectiva, el proyecto busca contribuir al desarrollo de metodologías analíticas que permitan aprovechar de
manera más integral la información disponible en el baloncesto profesional.
\cite{devlin2020identifyinggroupcontributionsnba}